Given two monoids $A$ and $B$, a monoid homomorphism $f : A \to B$ is a function such that $f(1_A) = 1_B$ and $f(xy) = f(x) f(y)$ for any $x, y \in A$.
We will treat $\N^+$ as a multiplicative monoid.

Let $M$ be an commutative and cancellative monoid.
Given a monoid homomorphism $f : \N^+ \to M$ and a positive integer $n$, we say that $n$ is $f$-\emph{nice} if the following holds:
  for any $a, b \in \N^+$ such that $a + b = n$, we have $f(a) = f(b)$.
Find all monoid homomorphism $f : \N^+ \to M$ such that there are infinitely many $f$-nice positive integers.



\subsection*{Answer}

The set of such monoid homomorphisms can be described as the union of several sets whose pairwise intersection is $\{0\}$.
We describe these smaller sets below.

First, fix a prime $p$, and consider the monoid $S_{p, M} = M \times \Hom(\F_p^{\times}/\{\pm 1\}, M)$.
We view $\Hom(\F_p^{\times}/\{\pm 1\}, M)$ as the set of monoid homomorphisms $\phi : \F_p^{\times} \to M$ such that $\phi(-1) = 1$.
For any pair $(c, \phi) \in S_{p, M}$, define the corresponding homomorphism $T_p(c, \phi) : \N^+ \to M$ by
\[ T_p(c, \phi)(p^k t) = c^k \phi(t) \]
  for any $k, t \in \N$ with $t$ coprime to $p$.
One can check that $T_p(c, \phi)$ is indeed a monoid homomorphism.
To see that $T_p(c, \phi)$ satisfies the main criterion, we claim that $p^k$ is $T_p(c, \phi)$-nice for any $k \geq 0$.
Indeed, for any $a, b \in \N^+$ with $a + b = p^k$, we can write $a = p^{\ell} a'$ and $b = p^{\ell} b'$ for some $\ell < k$ and $a', b' \in \N^+$ coprime with $p$.
But then $p \mid a' + b'$, so
\[ T_p(c, \phi)(a) = c^{\ell} \phi(a') = c^{\ell} \phi(b') = T_p(c, \phi)(b). \]
One can check that $T_p : S_{p, M} \to \Hom(\N^+, M)$ is an injective homomorphism.
We obtain the small set of solutions $T_p(S_{p, M})$ for each prime $p$.

Now, for the remaining class, consider an arbitrary sequence of pairs $(p_k, \phi_k)_{k \geq 1}$.
Here, $p_1 < p_2 < \ldots$ are primes and each $\phi_k : \Hom(\F_{p_k}^{\times}/\{\pm 1\}, M)$ are homomorphisms.
For a shorthand, we will denote this sequence by $(\phi_i)_{i \geq 1}$.
Let $S_{\infty, M}$ be the set of such sequences with the following additional property:
  for any $n, i, j \in \N^+$ with $n < \min\{p_i, p_j\}$, we have $\phi_i(n) = \phi_j(n)$.
Each sequence in $S_{\infty, M}$ gives rise to a monoid homomorphism $T_{\infty}((\phi_i)_{i \geq 1}) : \N^+ \to M$ given by
\[ T_{\infty}((\phi_i)_{i \geq 1})(n) = \phi_k(n) \]
  where $k$ is an arbitrary index with $n < p_k$.
Each $p_k$ is $T_{\infty}((\phi_i)_{i \geq 1})$-nice, so there are infinitely many $T_{\infty}((\phi_i)_{i \geq 1})$-nice primes.
The image $T_{\infty}(S_{\infty, M})$ is the remaining small set of solutions.
For each sequence $(\phi_i)_{i \geq 1}$, we have $T_{\infty}((\phi_i)_{i \geq 1}) = 1$ if and only if each $\phi_i$ is the constant one map.

In the case where $M$ has no torsion element, we get $\Hom(\F_p^{\times}/\{\pm 1\}, M) = \{1\}$ for each prime $p$.
Thus the set of solutions simplifies to homomorphisms of form $n \mapsto c^{\nu_p(n)}$ for some $p$ prime and $c \in M$.
Each $S_{p, M}$ is isomorphic to $M$, and $T_{\infty}(S_{\infty, M}) = \{1\}$.

Unfortunately, when $M$ has a torsion element, the set $T_{\infty}(S_{\infty, M})$ is non-trivial.
The case $M = \Z/2\Z$ can be solved using quadratic reciprocity and Dirichlet's theorem on arithmetic progressions.
The case $M = \Z/q\Z$ for arbitrary primes $q$ can be solved using generalizations: Eisenstein reciprocity and Chebotarev density theorem.
These cases can be combined to construct a non-trivial element of $T_{\infty}(S_{\infty, M})$ whenever $M$ is not torsion-free.



\subsection*{Solution}

Reference: \url{https://artofproblemsolving.com/community/c6h2625925p29340833}

We reference a solution in AoPS thread for IMO 2020 N5 by the user \textbf{BlazingMuddy} (see post \#40).
We have checked that each function in $T_p(S_{p, M})$ and $T_{\infty}(S_{\infty, M})$ works.
We now go to prove that they are the only functions that work.
We also prove that the $T_p(S_{p, M})$s and $T_{\infty}(S_{\infty, M})$ pairwise intersects at $\{1\}$.
We just run through several lemmas.

\begin{lemma}\label{2020n5-1}
Let $f : \N^+ \to M$ be a homomorphism.
Let $n \in \N^+$ be $f$-nice.
Then for any $d \in \N^+$ such that $d \mid n$, $d$ is $f$-nice.
\end{lemma}
\begin{proof}
Write $n = dk$ for some positive integer $k$.
Then for any $a, b \in \N^+$ such that $a + b = d$, we have $ak + bk = n$.
Since $n$ is $f$-nice,
\[ f(a) f(k) = f(ak) = f(bk) = f(b) f(k). \]
Thus $f(a) = f(b)$; this proves that $d$ is $f$-nice.
\end{proof}

\begin{lemma}\label{2020n5-2}
Let $f : \N^+ \to M$ be a homomorphism.
Let $p$ be a prime, and suppose that $p$ is $f$-nice.
Then $f(km \mod{p}) = f(k) f(m)$ for any positive integers $k, m < p$. 
Here $km \mod{p}$ denotes the remainder of $km$ under division by $p$.
\end{lemma}
\begin{proof}
Proceed by double strong induction on the order of $k$, then $m$.
The bases cases are all the cases with $km < p$; these are trivial cases.
We now describe the induction hypothesis and the goal to be proved.

Fix some $k_0, m_0 < p$, and WLOG assume that $k_0 m_0 > p$.
The induction hypothesis is that:
\begin{itemize}
    \item   for any $k < k_0$ and $m < p$, we have $f(km \mod{p}) = f(k) f(m)$; and
    \item   for any $m < m_0$, we have $f(k_0 m \mod{p}) = f(k_0) f(m)$.
\end{itemize}

The goal is to show that $f(k_0 m_0 \mod{p}) = f(k_0) f(m_0)$.
Write $p = q m_0 + r$ for some $q, r \in \N$.
Note that $0 < r < m_0$ and also $q < k_0$ since $p < k_0 m_0$.
Then by the second induction hypothesis and the fact that $p$ is $f$-nice, we get
\[ f(k_0 r \mod{p}) = f(k_0) f(r) = f(k_0) f(q m_0) = f(k_0) f(q) f(m_0). \]
But $p$ divides $k_0 r + k_0 q m_0 = k_0 p$.
Using the first induction hypothesis, we then get
\[ f(k_0 r \mod{p}) = f(k_0 q m_0 \mod{p}) = f(q) f(k_0 m_0 \mod{p}). \]
Finally, combine the two equalities and cancel out $f(q)$ from both sides.
Induction step is complete; the lemma is proved.
\end{proof}

The above lemma tells us that $f$ restricts to a homomorphism $\F_p^{\times} \to M$ if $p$ is $f$-nice.
Since we also require $f(p - 1) = f(1) = 1$, by the first isomorphism theorem, $f$ restricts to a homomorphism $\F_p^{\times}/\{\pm 1\} \to M$.

Due to Lemma~\ref{2020n5-1}, we have two cases regarding $f$-nice positive integers.
The first case is that there are infinitely many $f$-nice \emph{primes}.
In this case, enumerate the $f$-nice primes, say as $(p_k)_{k \geq 1}$.
By the argument on the previous paragraph, we get $f = T_{\infty}((\phi_i)_{i \geq 1})$ for some $(\phi_i)_{i \geq 1} \in S_{\infty, M}$.

The second, more structurally well-behaved case, is that there exists a prime $p$ such that $p^k$ is $f$-nice for every $k \geq 0$.
The following lemma allows us to prove that $f = T_p(c, \phi)$ for some $c \in M$ and $\phi : \F_p^{\times}/\{\pm 1\} \to M$.

\begin{lemma}\label{2020n5-3}
Let $f : \N^+ \to M$ be a homomorphism.
Let $p$ be a prime, and suppose that $p^k$ is $f$-nice for every $k \in \N$.
Then $f(n) = f(n \mod{p})$ for any $n \in \N^+$ coprime with $p$.
\end{lemma}
\begin{proof}
Induction on $n$; the base cases $n < p$ are easy.
Now suppose that $n > p$, and let $a$ be the positive integer such that $p^a < n < p^{a + 1}$.
Then we can write $p^{a + 1} = qn + r$ for some $q, r \in \N^+$ with $q < p$ and $r < n$.
Note that $r$ is coprime with $p$ since $q$ and $n$ are coprime with $p$.
By induction hypothesis, $f(r) = f(r \mod{p})$.
Since $p^{a + 1}$ and $p$ are $f$-nice (respectively), we get
\[ f(q) f(n) = f(qn) = f(r) = f(r \mod{p}) = f(qn \mod{p}). \]
By Lemma~\ref{2020n5-2}, since $q < p$, we get $f(qn \mod{p}) = f(q) f(n \mod{p})$.
Finally, cancel out $f(q)$ from both sides.
Induction step is complete.
\end{proof}

To finish from the lemma, recall by Lemma~\ref{2020n5-2} that $f$ restricts on $\{1, 2, \ldots, p - 1\}$ to a homomorphism $\phi : \F_p^{\times}/\{\pm 1\} \to M$.
Then $f(n) = \phi(n)$ for all $n$ coprime with $p$.
Thus we get $f = T_p(f(p), \phi)$.


\subsubsection*{Disjointness of the Small Sets of Solutions}

We now prove that the sets $T_p(S_{p, M})$ and $T_{\infty}(S_{\infty, M})$ only intersects at $\{1\}$, pairwise.
The following lemma suffices.

\begin{lemma}\label{2020n5-4}
Fix a prime $p$, and let $f \in T_p(S_{p, M})$.
Let $p$ be a prime and suppose that $p^k$ is $f$-nice for all $k \in \N$.
Suppose that there exists $n > p^2$ such that $n$ is $f$-nice.
Then $f \equiv 1$.
\end{lemma}
\begin{proof}
First, for any $0 < k \leq p$, we know that $f(kp) = f(n - kp)$.
By Lemma~\ref{2020n5-3}, $f(n - kp) = f((n - kp) \mod{p}) = f(n \mod{p})$ for all $k \leq p$.
In particular, this means $f(kp) = f(p)$, and thus $f(k) = 1$, for all $k \leq p$.

Now write $f = T_p(c, \rho)$ for some $c \in M$ and $\rho \in \Hom(\F_p^{\times}/\{\pm 1\}, M)$.
Then $c = f(p) = 1$ and $\rho(k) = f(k) = 1$ for all $0 < k < p$; the latter yields $\rho = 1$.
Thus $f = T_p(1, 1) = 1$, as desired.
\end{proof}

Now fix any $f \in T_p(S_{p, M})$.
First suppose that $f \in T_q(S_{q, M})$ for some prime $q \neq p$.
Without loss of generality, assume that $q > p$.
Then $f$ is $q^2$-good and $q^2 > p^2$, so the above lemma implies $f \equiv 1$.

Now suppose that $f = T_{\infty}((\phi_i)_{i \geq 1})$ for some $(\phi_i)_{i \geq 1} \in S_{\infty, M}$.
Consider the underlying sequence of primes $(p_i)_{i \geq 1}$.
We can choose $i$ large enough so that $p_i > p^2$, and $f$ would be $p_i$-good.
Then the above lemma implies $f \equiv 1$.


\subsubsection*{Non-triviality of $T_{\infty}(S_{\infty, M})$}

If $M$ is torsion-free and $N$ is a torsion monoid, then $\Hom(N, M) = 1$.
Thus Lemma~\ref{2020n5-2} yields that if $M$ is torsion-free and a prime $p$ is $f$-good, then $f(k) = 0$ for all $k < p$.
This reduces the solution's description to just $n \mapsto c^{\nu_p(n)}$ for some $c \in M$ and $p$ prime.
In particular, $T_p(S_{p, M}) \cong M$ and $T_{\infty}(S_{\infty, M}) = \{1\}$ in this case.

The case where $M$ is not torsion-free is harder to deal with.
We claim that actually, $T_{\infty}(S_{\infty, M})$ is non-trivial, i.e., not a singleton.
To prove the claim, it suffices to consider the cases where $M \cong C_q$, the cyclic group on $q$ elements, for some prime $q$.
We treat $C_q$ as a multiplicative group, isomorphic to $(\Z/q\Z, +)$.
We demonstrate just one case, $q = 2$.

For $q = 2$, consider $C_2 = \{\pm 1\}$.
The goal is to find a sequence $(p_i, \phi_i)_{i \geq 1}$ such that $\phi_i \neq 1$ for some $i$ and also $\phi_i(n) = \phi_j(n)$ for any $n < p_i, p_j$.
We can take each $\phi_i$ to be the Legendre symbol $\left(\frac{\cdot}{p_i}\right)$.
The criterion $\phi_i(-1) = 1$ amounts to $p_i \equiv 1 \pmod{4}$ for all $i \geq 1$.
Due to quadratic reciprocity, the condition $\phi_i(n) = \phi_j(n)$ for any $n < p_i, p_j$ is fulfilled whenever $p_{i + 1} \equiv p_i \pmod{4(p_i - 1)!}$ for all $i = 1, 2, \ldots$.
Such sequence of primes exist due to Dirichlet's theorem on arithmetic progressions.

For the general case on $q$, one has to work with the $q^{\text{th}}$ power symbol over $\Z[\zeta_q]$.
The reciprocity result we need is called Eisenstein reciprocity.
Construction of the desired primes/prime ideals follow from a corollary of Chebotarev's density theorem.



\subsection*{Implementation details}

The results we have to prove have disjointness from each other to some degree.
They will be split into several files.

Unfortunately, \texttt{mathlib4} has a way better API on \texttt{Nat} than \texttt{PNat}.
Also, it seems to me that using $\Hom(\N^+, M)$ in the implementation results in slower definitions and lemmas.
Thus, we implement monoid homomorphisms as a function $f : \N \to M$ satisfying $f(0) = f(1) = 1$ and $f(xy) = f(x) f(y)$ for $x, y \neq 0$.
